\documentclass[11pt]{article}
\usepackage[a4paper,margin=1in]{geometry}
\usepackage{fontspec}
\usepackage[boldfont=false]{xeCJK}
\setCJKmainfont{FandolSong-Regular.otf}
\usepackage{amsmath}
\usepackage{amssymb}
\usepackage{array}
\usepackage{booktabs}
\usepackage{graphicx}
\usepackage{adjustbox}
\usepackage{enumitem}
\setlist[itemize]{topsep=2pt,partopsep=0pt,parsep=0pt,itemsep=2pt,leftmargin=1.4em}
\setlist[enumerate]{topsep=2pt,partopsep=0pt,parsep=0pt,itemsep=2pt,leftmargin=1.6em}
\usepackage{parskip}
\usepackage{fvextra}
\fvset{breaklines=true,breakanywhere=true,fontsize=\small}
\setlength{\parindent}{0pt}

\begin{document}

\begin{center}
  {\LARGE\bfseries Marketing}\\[0.35em]
  {\large A Level Business}
\end{center}
\vspace{0.6em}

\section*{The role of marketing}

\textbf{Marketing} 市场营销 is the job of finding out what customers want, then making and selling products that meet those wants at a profit. Good marketing links the business to its \textbf{market} 市场 (all the buyers and sellers of a product).

A business can be:

\begin{itemize}
\item \textbf{product-orientated} 产品导向 — it focuses on the product first and hopes customers will buy it. This suits new or technical products.

\item \textbf{market-orientated} 市场导向 — it finds out what customers want first, then makes it. This lowers the risk of failure.

\end{itemize}

\textbf{Customer relationship marketing} 客户关系营销 (CRM) means building long-term links with customers so they come back again. Keeping a customer is cheaper than finding a new one.

\section*{Market size, growth and share}

\begin{itemize}
\item \textbf{market size} 市场规模 — the total sales of all firms in a market, by value or by number of units.

\item \textbf{market growth} 市场增长 — how fast the market is getting bigger, shown as a percentage.

\item \textbf{market share} 市场份额 — one firm's sales as a percentage of the total market:

\end{itemize}

\[
\text{market share} = \frac{\text{firm's sales}}{\text{total market sales}} \times 100\%
\]

A rising market share is a sign of strong marketing.

\section*{Mass and niche markets}

\begin{itemize}
\item a \textbf{mass market} 大众市场 is large, with products aimed at most buyers (e.g. soft drinks). Sales are high, but competition is strong.

\item a \textbf{niche market} 利基市场 is a small part of a market with special needs (e.g. left-handed tools). Competition is weaker, but sales are smaller.

\end{itemize}

\section*{Market segmentation}

\textbf{Market segmentation} 市场细分 means splitting a market into groups (\textbf{segments} 细分市场) of buyers who are alike. Common ways to segment are by age, income, gender, location, and lifestyle. Segmentation lets a firm aim the right product and message at each group, which raises sales and cuts waste.

\section*{Demand, supply and elasticity}

\textbf{Demand} 需求 is how much of a product buyers will buy at each price. \textbf{Supply} 供给 is how much sellers will offer at each price. Usually, when price falls, demand rises.

\textbf{Elasticity} 弹性 measures how strongly demand reacts to a change. You should know three types.

\textbf{Price elasticity of demand} 需求价格弹性 (PED) measures how demand reacts to a price change:

\[
\text{PED} = \frac{\%\ \text{change in quantity demanded}}{\%\ \text{change in price}}
\]

\begin{itemize}
\item if PED is greater than 1, demand is \textbf{price elastic} 富有弹性 — a small price cut raises sales a lot, so lowering price can raise revenue.

\item if PED is less than 1, demand is \textbf{price inelastic} 缺乏弹性 — demand changes little, so raising price can raise revenue.

\end{itemize}

\textbf{Income elasticity of demand} 需求收入弹性 (YED) measures how demand reacts to a change in customers' income. Luxury goods react strongly; basic goods react weakly.

\textbf{Cross elasticity of demand} 需求交叉弹性 (XED) measures how the demand for one product reacts to a price change of another. It is positive for \textbf{substitutes} 替代品 (e.g. tea and coffee) and negative for \textbf{complements} 互补品 (e.g. cars and fuel).

Elasticity is useful because it helps a firm predict how a price change will affect sales and revenue.

\section*{Market research}

\textbf{Market research} 市场调研 is collecting and studying information about customers, \textbf{competitors} 竞争对手 and the market.

\begin{itemize}
\item \textbf{primary research} 一手调研 — new information you collect yourself, by surveys, interviews or observation. It fits your exact needs, but costs more and takes time.

\item \textbf{secondary research} 二手调研 — information that already exists, such as reports, websites and sales records. It is cheap and quick, but may be old or not exactly right.

\end{itemize}

Data can be:

\begin{itemize}
\item \textbf{quantitative data} 定量数据 — facts in numbers (e.g. "60\% buy weekly"). Easy to compare.

\item \textbf{qualitative data} 定性数据 — opinions and reasons (e.g. "why people like the brand"). Gives deeper understanding.

\end{itemize}

\subsection*{Sampling and reliability}

A \textbf{sample} 样本 is the small group of people asked, chosen to stand for the whole market. \textbf{Sampling} 抽样 methods include:

\begin{itemize}
\item \textbf{random sampling} 随机抽样 — everyone has an equal chance of being picked.

\item \textbf{quota sampling} 配额抽样 — set numbers from each group are asked.

\item \textbf{stratified sampling} 分层抽样 — the sample copies the make-up of the whole market.

\end{itemize}

A larger, well-chosen sample gives more \textbf{reliable} 可靠的 results. Poor sampling can cause \textbf{bias} 偏差, which makes the findings wrong. Results are often shown in tables and charts, then read carefully to guide decisions.

\section*{The marketing mix}

The \textbf{marketing mix} 营销组合 is the set of choices a firm makes to sell a product. It is often called the \textbf{4Ps}:

\begin{center}
\adjustbox{max width=\linewidth}{%
\begin{tabular}{ll}
\toprule
\textbf{P} & \textbf{Question it answers} \\
\midrule
\textbf{product} 产品 & what are we selling? \\
\textbf{price} 价格 & how much do we charge? \\
\textbf{place} 渠道 & where and how do customers buy it? \\
\textbf{promotion} 促销 & how do we tell customers about it? \\
\bottomrule
\end{tabular}}
\end{center}

The four parts must work together. For example, a high price needs high quality and a good image. Making the parts fit each other is called \textbf{integration} 整合.

\section*{Product}

A good product is the heart of the mix.

\textbf{Product development} 产品开发 is the work of designing and improving products to meet customer needs and beat rivals.

\subsection*{The product life cycle}

The \textbf{product life cycle} 产品生命周期 shows the sales of a product over time, in four stages:

\begin{enumerate}
\item \textbf{introduction} 引入期 — launch; low sales, high promotion cost.

\item \textbf{growth} 成长期 — sales rise quickly.

\item \textbf{maturity} 成熟期 — sales reach their peak and level off.

\item \textbf{decline} 衰退期 — sales fall as tastes change.

\end{enumerate}

To keep sales up in maturity, firms use \textbf{extension strategies} 延长策略, such as new packaging, new uses, or selling in new markets.

\subsection*{The product portfolio (Boston Matrix)}

The \textbf{product portfolio} 产品组合 is the full range of products a firm sells. The \textbf{Boston Matrix} 波士顿矩阵 sorts products by market share and market growth:

\begin{center}
\adjustbox{max width=\linewidth}{%
\begin{tabular}{lll}
\toprule
\textbf{} & \textbf{High market growth} & \textbf{Low market growth} \\
\midrule
\textbf{High share} & \textbf{stars} 明星产品 & \textbf{cash cows} 现金牛 \\
\textbf{Low share} & \textbf{question marks} 问题产品 & \textbf{dogs} 瘦狗产品 \\
\bottomrule
\end{tabular}}
\end{center}

Stars need investment but may become cash cows. Cash cows bring steady profit. Question marks are risky. Dogs are usually dropped.

\subsection*{Branding}

\textbf{Branding} 品牌 gives a product a name, design and image that customers recognise. A strong brand builds \textbf{loyalty} 忠诚度, lets the firm charge more, and makes new products easier to launch.

\section*{Price}

Choosing a price is a key decision. The main methods and strategies are:

\begin{center}
\adjustbox{max width=\linewidth}{%
\begin{tabular}{ll}
\toprule
\textbf{Method / strategy} & \textbf{How it works} \\
\midrule
\textbf{cost-plus pricing} 成本加成定价 & add a fixed profit margin to the unit cost \\
\textbf{penetration pricing} 渗透定价 & set a low price to enter a market and win share \\
\textbf{price skimming} 撇脂定价 & set a high price first for a new, special product, then lower it \\
\textbf{competitive pricing} 竞争定价 & set a price close to rivals' prices \\
\textbf{price discrimination} 价格歧视 & charge different prices to different groups (e.g. peak vs off-peak) \\
\textbf{dynamic pricing} 动态定价 & change the price often as demand changes \\
\textbf{psychological pricing} 心理定价 & use prices like \$9.99 that feel lower \\
\bottomrule
\end{tabular}}
\end{center}

The best price depends on costs, competitors, the product's life-cycle stage, and how price elastic demand is.

\section*{Place and promotion}

\textbf{Place} is how the product reaches the customer, through a \textbf{distribution channel} 分销渠道. A short channel (maker → customer) gives more control; a longer channel (maker → \textbf{wholesaler} 批发商 → shop → customer) reaches more buyers.

\textbf{Promotion} tells customers about a product and persuades them to buy. It is split into:

\begin{itemize}
\item \textbf{above-the-line promotion} 大众媒体广告 — paid \textbf{advertising} 广告 in mass media, such as TV and online ads, aimed at a wide audience.

\item \textbf{below-the-line promotion} 直接促销 — targeted methods the firm controls, such as discounts, free samples and direct mail.

\end{itemize}

\textbf{Digital marketing} 数字营销 and \textbf{e-commerce} 电子商务 (buying and selling online) now play a huge role. They let small firms reach world markets cheaply, target adverts to the right people, and sell at any time of day.


\end{document}
